\scanpage{11}
\uline{Théorème}\qquad (i) (Pour mémoire) $p\colon\mathscr B\longrightarrow G'$ est projectif,
\[
\widetilde f\colon\mathscr B\longrightarrow T\quad\text{est lisse.}
\]

(ii) $p$ induit
\[
\boxed{\mathscr B^{\mathrm{reg}}\overset{\mathrm{dfn}}{=}\mathscr B^{\deleted{\unsure{\mathrm{reg}}}}|G'^{\mathrm{reg}}
\xrightarrow{\sim}G'^{\mathrm{reg}}}
\]
et $G'^{\mathrm{reg}}$ est le plus grand ouvert au dessus duquel $p$ est un isom.

(iii) (Pour mémoire) $\mathscr B-\mathscr B^{\mathrm{reg}}$ est de codim $\geq 2$ \deleted{\unsure{$\mathscr B$}} dans $\mathscr B$ (fibre par fibres sur $S$ (ou sur $D$)

$G'-G'^{\mathrm{reg}}$ a toutes ses composantes irréductibles de codimension 3 (fibre par fibre sur $S$)

(iv) \inserted{(Si $S$ \deleted{\unsure{régulier}}\inserted{normal} \deleted{\illegible{1}} loc. noeth.)}
\diagram{\begin{tikzcd}[ampersand replacement=\&,column sep=small]
\mathrm{Pic}(\mathscr B)\arrow[r,"\sim"] \& \mathrm{Pic}(\mathscr B^{\mathrm{reg}}) \& \mathrm{Pic}(G'^{\mathrm{reg}})\arrow[l,"\sim"']\ \deleted{\unsure{\mathrm{Pic}(G')}}\\
\mathrm{Pic}(D)\arrow[u,"{\wr}"'] \& \&
\end{tikzcd}}
\inserted{[Si $T$ non tordu, \deleted{\unsure{ceci signifie simplement}} \unsure{resp. plus intrinsèquement|ou plus intrinsèquement} \unsure{généralement|génériquement} si pour tt point maximal $\unsure{\eta|s|\xi}$ de $S$, l'image $\Gamma$ de $\pi_1(S,\bar\eta)$ dans $\operatorname{Ext}(G_{\bar\eta})$ est telle que $H^1(\Gamma,M)=0$] on a une réciproque si $S$ régulier.}

et en faisceautisant sur $S$ (pour la top. \unsure{éte|ét.|étale}, disons)
\diagram{\begin{tikzcd}[ampersand replacement=\&,column sep=small]
\& \underline{\mathrm{Pic}}_{\mathscr B/S}\arrow[r,"\sim"] \& \underline{\mathrm{Pic}}_{\mathscr B^{\mathrm{reg}}/S} \& \underline{\mathrm{Pic}}_{G'^{\mathrm{reg}}/S}\arrow[l,"\sim"']\\
M\arrow[r]\&P\simeq\underline{\mathrm{Pic}}_{D/S}\arrow[u,"{\wr}"',"{\text{\deleted{\unsure{réciproque}}}}" near start,"{\text{\deleted{\unsure{le même groupe}}}}"' near end] \& \&
\end{tikzcd}}
\inserted{(injectif si $G$ semisimple)}

De plus, le cône ample pour $\underline{\mathrm{Pic}}_{\mathscr B/S}$ est le $=$ que pour $\underline{\mathrm{Pic}}_{D/S}$, qui est le cône des poids $P^+$. \deleted{\unsure{Et on a}}

\uline{Démonstration}\quad (i) et (iii) connus (pour (iii) ref : Steinberg ....)

(ii) On commence à voir, posant \deleted{$G^{rs}=\unsure{\text{ouv}}$}

$G^{rs}\subset G^{\mathrm{reg}}$ ouvert des pts de $G$ qui sont réguliers semi-simples
\[
\left.\begin{aligned}
G'^{rs}&\subset G^{\mathrm{reg}}\\
\mathscr B^{rs}&\subset\mathscr B^{\mathrm{reg}}
\end{aligned}\right|\quad\text{images inverses,}
\]
que
\[
\mathscr B^{rs}\xrightarrow{\sim} G'^{rs}
\]
D'autre part, on sait par Steinberg que les fibres de $\mathscr B^{\mathrm{reg}}\longrightarrow G'^{\mathrm{reg}}$ sont finies, d'ailleurs c'est un morphisme propre. Comme $G'^{\mathrm{reg}}$ est normal, on conclut par le Main Theorem.
\margin{On sait aussi par Steinberg que les fibres des autres pts sont de dim $>0$.}
% archive sheet number: 10
